%% rset.rigor.tex
%% What it would take to turn the programmatic construction into a rigorous one.
%% Drop in via %% rset.rigor.tex
%% What it would take to turn the programmatic construction into a rigorous one.
%% Drop in via %% rset.rigor.tex
%% What it would take to turn the programmatic construction into a rigorous one.
%% Drop in via %% rset.rigor.tex
%% What it would take to turn the programmatic construction into a rigorous one.
%% Drop in via \input{rset.rigor} (after rset.nwf, before rset.future).
\input{rset.colors}   % red/black colour layer + FM symbols (providecommand-guarded)

\section{Toward a rigorous construction}
\label{sec:rigor}

These are notes, and the constructions above are pitched at the level of intent:
they say what the theory is meant to be and exhibit the model it is meant to
carry, but several of the moves are made by notation where they ought to be made
by theorem. We close by setting down, as a checklist of definitions to fix and
obligations to discharge, what a fully rigorous development would require. We
group them under the headings a careful referee would use --- the fixpoint, the
metatheory, the membership theory, comprehension, and the construction's place
among its neighbours --- and note that for the finitary core each obligation
reduces to a statement about hereditarily finite sets.

\subsection{The knotted fixpoint, via algebraic set theory}
\label{sec:rigor:fixpoint}

The central move, $\rX = \bSetX$ and $\bX = \rSetX$, is at present a definition by
notation; it must be exhibited as a fixed point that provably exists. Algebraic
set theory ($\mathsf{AST}$) is, we think, the right basis for this, because it is
built precisely to produce universes of sets as fixed points of a
power\emph{class} operator while dodging the size obstruction that defeats the
naive power\emph{set} construction \cite{joyal1995algebraic,awodey2008ast}.

Recall its shape. One works in a Heyting category $\mathcal{C}$ of \emph{classes}
with a distinguished collection of \emph{small} maps and a \emph{powerclass}
functor $\mathcal{P}_s\colon\mathcal{C}\to\mathcal{C}$, sending a class to the
class of its \emph{small} sub-objects. The universe of sets is the free
$\mathcal{P}_s$-algebra: an object $U$ with a structural isomorphism
$\mathcal{P}_s U\cong U$, and the central theorem is that this initial
$\mathsf{ZF}$-algebra exists in any category with small maps and models the axioms
of set theory, with membership recovered from the algebra structure. Restricting
$\mathcal{P}_s$ to \emph{small} sub-objects, rather than arbitrary sub-classes, is
exactly what makes $\mathcal{P}_s U\cong U$ possible where Cantor forbids
$\mathcal{P}X\cong X$ --- so the size problem flagged in Section~\ref{sec:nwf} is
dissolved at the level of the framework, not patched case by case.

Two features of $\mathsf{AST}$ fit the present construction almost too well.
First, it handles atoms by a coproduct: the cumulative hierarchy over an object
$A$ of atoms is the initial algebra of $Y\mapsto A + \mathcal{P}_s Y$, in which the
points of $A$ acquire no members --- they are atoms because the free construction
\emph{forgets} whatever structure they carried. That forgetfulness is precisely
our colour opacity (Remark~\ref{rem:group}): red sees the points of its
atom-supply but not their black interior. Second, the universe is constructed as
a $W$-type quotiented by a bounded bisimulation, so the development is
constructive and portable --- and, not incidentally, close to what a
formalisation would have to build.

The knot is then a two-sorted instance. Take the endofunctor on
$\mathcal{C}\times\mathcal{C}$
\[
  H(R,B) \;=\; \bigl(\,B + \mathcal{P}_s R,\ \ R + \mathcal{P}_s B\,\bigr),
\]
whose initial algebra is a pair $(\rX,\bX)$ with $\rX\cong\bX + \mathcal{P}_s\rX$
and $\bX\cong\rX + \mathcal{P}_s\bX$: each colour's universe is its own sets
together with the other colour's universe \emph{as atoms} --- precisely
$\rX=\bSetX$ and $\bX=\rSetX$. The obligation thus becomes to run the
$\mathsf{AST}$ existence and modelling theorems for $H$ in place of
$\mathcal{P}_s$ alone. The ingredients are unchanged --- coproduct, powerclass,
$W$-types, bounded bisimulation --- so one expects the same proof to go through
two-sorted; what it would establish is that the knotted universe \emph{exists} and
that each colour \emph{models} $\mathsf{ZFA}$, discharging in a single stroke both
this obligation and the extensionality and separation obligations of
Section~\ref{sec:rigor:ext}.

Three riders complete the picture. For the \emph{finitary} core one replaces
$\mathcal{P}_s$ by the finite-powerclass functor; being finitary, the initial
algebra of the corresponding $H$ exists already at the level of \emph{sets}, as
the colimit of the initial $\omega$-chain $\varnothing\to H(\varnothing)\to
H^2(\varnothing)\to\cdots$, with no class machinery required --- and it is the
$\mathrm{HF}$ universe of Section~\ref{sec:fmmodel}. For the non-well-founded
reading of Section~\ref{sec:nwf} one takes the \emph{final coalgebra} of $H$ in
place of the initial algebra: at set size for the finitary functor, and in the
class framework --- where it coincides with Aczel's universe --- for the full one.
And a caveat worth stating: the initial $\mathsf{ZF}$-algebra models
\emph{constructive} set theory by default, so recovering classical
$\mathsf{ZFA}$/$\mathsf{FM}$ requires the corresponding classical principles
(power set, full separation, excluded middle) on the class category. With those in
place, $\mathsf{AST}$ supplies exactly the existence theorem these notes have so
far only gestured at, and supplies it in a form that already carries the
membership and extensionality structure the rest of this section asks for.

\subsection{Metatheory and consistency strength}
\label{sec:rigor:meta}

The notes never name the theory in which they are conducted, and the headline
claim --- atoms without infinite risk --- has content only relative to a fixed
metatheory. The obligation is to state that metatheory and to cash the risk claim
as a statement of interpretability. The natural target is hereditarily finite set
theory, equivalently first-order arithmetic: the model $\mathcal{M}$ of
Theorem~\ref{thm:rep} is itself an interpretation of the finitary reflective
theory in $\mathrm{HF}$, so the two are equiconsistent and the ``risk'' is exactly
that of arithmetic --- no infinite set of primitive atoms is ever posited. The
representation theorem, in other words, already \emph{is} a relative-consistency
proof; what remains is to state it as one. The infinitary extension of
Section~\ref{sec:fmmodel:infinitary} raises the target by a single $\omega$: it is
interpretable in $\mathsf{Z}$ --- or in $\mathsf{ZF}$ with foundation reinstated
per colour --- the lone infinity being black infinity. Making ``finite versus
infinite risk'' precise in this way is, we expect, the first thing a set theorist
would ask for, and it is the claim on which the paper's title rests.

\subsection{Extensionality and the membership theory}
\label{sec:rigor:ext}

Extensionality is not visibly secured by the equations of the formal theory, which
give reflexivity, symmetry, transitivity, and idempotence of $\rcup$ and $\rcap$,
but not that sets with the same members are equal. A rigorous development must
either adopt extensionality as an axiom or derive it, and must then prove that
provable red equality coincides with sameness of members --- the property the
proof of Theorem~\ref{thm:rep} leans on. The membership rules need a matching
completeness theorem. As presented they decide $\rin$ on unions, intersections,
differences, and singletons, but no rule governs membership in a nested
$\rbr{S}$, and one wants a proof that for $\mathrm{HF}$ sets every true $e\rin S$
and $e\rnin S$ is derivable and that the relation is decidable. Until then, the
remark that the syntactic theory is ``too fine grained'' and is corrected by
equations is a programme rather than a theorem. We note that the
algebraic-set-theory route of Section~\ref{sec:rigor:fixpoint} would discharge
this obligation along with existence: its modelling theorem delivers
extensionality for the constructed universe directly, so on that route the
equational theory need only be shown \emph{sound} for the model rather than
independently complete.

\subsection{Comprehension and freedom from paradox}
\label{sec:rigor:compr}

The $\mathrm{Comprehension}$ rule forms $\rbr{f(s_1,\dots,s_n) : s_i \rin S_i}$ for
a function $f\colon S_1\times\dots\times S_n\to S$, and the admissible class of
$f$ is left open although it governs both the strength of the theory and its
safety. The rule is the image of a function over sets already built ---
separation- and replacement-flavoured rather than naive comprehension --- so the
expected outcome is that no Russell set arises; but this must be argued, with
care in the cross-colour case, where $f$ must respect the opacity of atoms. The
cleanest argument is semantic and is already to hand: the $\mathrm{HF}$ model
witnesses consistency, so fixing $f$ to range over the metatheory's set functions
and reading the rule in $\mathcal{M}$ shows the theory paradox-free relative to
$\mathrm{HF}$. Specifying $f$ precisely, and confirming that the cross-colour
instances preserve well-formedness, is the remaining work.

\subsection{Placing the construction}
\label{sec:rigor:place}

Finally, the construction should be located among its neighbours, which would
both guard against reinvention and sharpen its claims. Three comparisons are
immediate. The permutation action and the atoms align it with the
Fraenkel--Mostowski permutation models and their nominal descendants
\cite{DBLP:journals/fac/GabbayP02}; the non-well-founded reading aligns it with
Aczel \cite{aczel1988nwf}; and the colour barrier --- red cannot see into black
--- is a two-level \emph{stratification}, inviting comparison with the typed and
stratified set theories. Most pointedly, red and black \emph{mutually interpret}
one another, so the construction is an instance of mutual interpretation between
set theories; whether the two copies are \emph{bi-interpretable}, and in what
precise sense the universe of set theories promised by the abstract is a
\emph{multiverse}, are exactly the questions that would connect these notes to the
contemporary study of interpretations between set theories. We regard settling
them as the natural completion of the present work.

None of these obligations looks out of reach for the finitary core, where each
reduces to a statement about hereditarily finite sets and several are corollaries
of the model already built. Together they are what would carry the construction
from a suggestive picture to a theorem --- and, not incidentally, they are close
to the list of lemmas a formalisation would have to state in any case.
 (after rset.nwf, before rset.future).
%% rset.colors.tex
%% Red/black colour layer for the two intertwined copies, plus the FM symbols
%% used by rset.mainthm.tex and rset.applications.tex.
%%
%% EVERYTHING routes through \rd (red copy = player) and \bk (black copy =
%% opponent): redefine just those two macros to change the whole scheme (e.g.
%% to make the "black" copy dark blue for legibility).  Every definition is
%% \providecommand-guarded, so it is safe to \input this file from several
%% places and harmless if rset.local already defines these names.
%%
%% Requires \usepackage{xcolor} (already loaded in rset.tex).

% --- the two copies -------------------------------------------------------
\providecommand{\rd}[1]{\textcolor{red}{#1}}     % red copy  (player)
\providecommand{\bk}[1]{\textcolor{black}{#1}}   % black copy (opponent)

% --- theory names and atom supplies --------------------------------------
\providecommand{\rSetX}{\rd{\mathrm{Set}[X]}}    % red   Set[X]
\providecommand{\bSetX}{\bk{\mathrm{Set}[X]}}    % black Set[X]
\providecommand{\rX}{\rd{X}}                      % red   atom supply
\providecommand{\bX}{\bk{X}}                      % black atom supply

% --- empty set ------------------------------------------------------------
\providecommand{\remp}{\rd{\emptyset}}
\providecommand{\bemp}{\bk{\emptyset}}

% --- membership (relations) ----------------------------------------------
\providecommand{\rin}{\mathrel{\rd{\in}}}
\providecommand{\bin}{\mathrel{\bk{\in}}}
\providecommand{\rnin}{\mathrel{\rd{\notin}}}
\providecommand{\bnin}{\mathrel{\bk{\notin}}}

% --- binary operations ----------------------------------------------------
\providecommand{\rcup}{\mathbin{\rd{\cup}}}
\providecommand{\bcup}{\mathbin{\bk{\cup}}}
\providecommand{\rcap}{\mathbin{\rd{\cap}}}
\providecommand{\bcap}{\mathbin{\bk{\cap}}}
\providecommand{\rsm}{\mathbin{\rd{\setminus}}}
\providecommand{\bsm}{\mathbin{\bk{\setminus}}}

% --- comprehension braces of each theory ---------------------------------
\providecommand{\rbr}[1]{\rd{\{}\,#1\,\rd{\}}}
\providecommand{\bbr}[1]{\bk{\{}\,#1\,\bk{\}}}

% --- colour-blind membership (sees through the colour barrier) ------------
\providecommand{\cbin}{\mathrel{\in_{\!\star}}}

% --- FM symbols (shared with rset.stdfmset.tex) --------------------------
\providecommand{\Atoms}{\mathbb{A}}
\providecommand{\Sym}{\mathrm{Sym}}
\providecommand{\pow}{\mathcal{P}}
\providecommand{\supp}{\mathrm{supp}}
\providecommand{\fs}{\mathrm{fs}}
\providecommand{\new}{\mathbin{\reflectbox{\ensuremath{\mathsf{N}}}}}
\providecommand{\atomsof}{\mathrm{atoms}}
\providecommand{\Red}{\rd{\mathsf{R}}}           % the red HF universe
   % red/black colour layer + FM symbols (providecommand-guarded)

\section{Toward a rigorous construction}
\label{sec:rigor}

These are notes, and the constructions above are pitched at the level of intent:
they say what the theory is meant to be and exhibit the model it is meant to
carry, but several of the moves are made by notation where they ought to be made
by theorem. We close by setting down, as a checklist of definitions to fix and
obligations to discharge, what a fully rigorous development would require. We
group them under the headings a careful referee would use --- the fixpoint, the
metatheory, the membership theory, comprehension, and the construction's place
among its neighbours --- and note that for the finitary core each obligation
reduces to a statement about hereditarily finite sets.

\subsection{The knotted fixpoint, via algebraic set theory}
\label{sec:rigor:fixpoint}

The central move, $\rX = \bSetX$ and $\bX = \rSetX$, is at present a definition by
notation; it must be exhibited as a fixed point that provably exists. Algebraic
set theory ($\mathsf{AST}$) is, we think, the right basis for this, because it is
built precisely to produce universes of sets as fixed points of a
power\emph{class} operator while dodging the size obstruction that defeats the
naive power\emph{set} construction \cite{joyal1995algebraic,awodey2008ast}.

Recall its shape. One works in a Heyting category $\mathcal{C}$ of \emph{classes}
with a distinguished collection of \emph{small} maps and a \emph{powerclass}
functor $\mathcal{P}_s\colon\mathcal{C}\to\mathcal{C}$, sending a class to the
class of its \emph{small} sub-objects. The universe of sets is the free
$\mathcal{P}_s$-algebra: an object $U$ with a structural isomorphism
$\mathcal{P}_s U\cong U$, and the central theorem is that this initial
$\mathsf{ZF}$-algebra exists in any category with small maps and models the axioms
of set theory, with membership recovered from the algebra structure. Restricting
$\mathcal{P}_s$ to \emph{small} sub-objects, rather than arbitrary sub-classes, is
exactly what makes $\mathcal{P}_s U\cong U$ possible where Cantor forbids
$\mathcal{P}X\cong X$ --- so the size problem flagged in Section~\ref{sec:nwf} is
dissolved at the level of the framework, not patched case by case.

Two features of $\mathsf{AST}$ fit the present construction almost too well.
First, it handles atoms by a coproduct: the cumulative hierarchy over an object
$A$ of atoms is the initial algebra of $Y\mapsto A + \mathcal{P}_s Y$, in which the
points of $A$ acquire no members --- they are atoms because the free construction
\emph{forgets} whatever structure they carried. That forgetfulness is precisely
our colour opacity (Remark~\ref{rem:group}): red sees the points of its
atom-supply but not their black interior. Second, the universe is constructed as
a $W$-type quotiented by a bounded bisimulation, so the development is
constructive and portable --- and, not incidentally, close to what a
formalisation would have to build.

The knot is then a two-sorted instance. Take the endofunctor on
$\mathcal{C}\times\mathcal{C}$
\[
  H(R,B) \;=\; \bigl(\,B + \mathcal{P}_s R,\ \ R + \mathcal{P}_s B\,\bigr),
\]
whose initial algebra is a pair $(\rX,\bX)$ with $\rX\cong\bX + \mathcal{P}_s\rX$
and $\bX\cong\rX + \mathcal{P}_s\bX$: each colour's universe is its own sets
together with the other colour's universe \emph{as atoms} --- precisely
$\rX=\bSetX$ and $\bX=\rSetX$. The obligation thus becomes to run the
$\mathsf{AST}$ existence and modelling theorems for $H$ in place of
$\mathcal{P}_s$ alone. The ingredients are unchanged --- coproduct, powerclass,
$W$-types, bounded bisimulation --- so one expects the same proof to go through
two-sorted; what it would establish is that the knotted universe \emph{exists} and
that each colour \emph{models} $\mathsf{ZFA}$, discharging in a single stroke both
this obligation and the extensionality and separation obligations of
Section~\ref{sec:rigor:ext}.

Three riders complete the picture. For the \emph{finitary} core one replaces
$\mathcal{P}_s$ by the finite-powerclass functor; being finitary, the initial
algebra of the corresponding $H$ exists already at the level of \emph{sets}, as
the colimit of the initial $\omega$-chain $\varnothing\to H(\varnothing)\to
H^2(\varnothing)\to\cdots$, with no class machinery required --- and it is the
$\mathrm{HF}$ universe of Section~\ref{sec:fmmodel}. For the non-well-founded
reading of Section~\ref{sec:nwf} one takes the \emph{final coalgebra} of $H$ in
place of the initial algebra: at set size for the finitary functor, and in the
class framework --- where it coincides with Aczel's universe --- for the full one.
And a caveat worth stating: the initial $\mathsf{ZF}$-algebra models
\emph{constructive} set theory by default, so recovering classical
$\mathsf{ZFA}$/$\mathsf{FM}$ requires the corresponding classical principles
(power set, full separation, excluded middle) on the class category. With those in
place, $\mathsf{AST}$ supplies exactly the existence theorem these notes have so
far only gestured at, and supplies it in a form that already carries the
membership and extensionality structure the rest of this section asks for.

\subsection{Metatheory and consistency strength}
\label{sec:rigor:meta}

The notes never name the theory in which they are conducted, and the headline
claim --- atoms without infinite risk --- has content only relative to a fixed
metatheory. The obligation is to state that metatheory and to cash the risk claim
as a statement of interpretability. The natural target is hereditarily finite set
theory, equivalently first-order arithmetic: the model $\mathcal{M}$ of
Theorem~\ref{thm:rep} is itself an interpretation of the finitary reflective
theory in $\mathrm{HF}$, so the two are equiconsistent and the ``risk'' is exactly
that of arithmetic --- no infinite set of primitive atoms is ever posited. The
representation theorem, in other words, already \emph{is} a relative-consistency
proof; what remains is to state it as one. The infinitary extension of
Section~\ref{sec:fmmodel:infinitary} raises the target by a single $\omega$: it is
interpretable in $\mathsf{Z}$ --- or in $\mathsf{ZF}$ with foundation reinstated
per colour --- the lone infinity being black infinity. Making ``finite versus
infinite risk'' precise in this way is, we expect, the first thing a set theorist
would ask for, and it is the claim on which the paper's title rests.

\subsection{Extensionality and the membership theory}
\label{sec:rigor:ext}

Extensionality is not visibly secured by the equations of the formal theory, which
give reflexivity, symmetry, transitivity, and idempotence of $\rcup$ and $\rcap$,
but not that sets with the same members are equal. A rigorous development must
either adopt extensionality as an axiom or derive it, and must then prove that
provable red equality coincides with sameness of members --- the property the
proof of Theorem~\ref{thm:rep} leans on. The membership rules need a matching
completeness theorem. As presented they decide $\rin$ on unions, intersections,
differences, and singletons, but no rule governs membership in a nested
$\rbr{S}$, and one wants a proof that for $\mathrm{HF}$ sets every true $e\rin S$
and $e\rnin S$ is derivable and that the relation is decidable. Until then, the
remark that the syntactic theory is ``too fine grained'' and is corrected by
equations is a programme rather than a theorem. We note that the
algebraic-set-theory route of Section~\ref{sec:rigor:fixpoint} would discharge
this obligation along with existence: its modelling theorem delivers
extensionality for the constructed universe directly, so on that route the
equational theory need only be shown \emph{sound} for the model rather than
independently complete.

\subsection{Comprehension and freedom from paradox}
\label{sec:rigor:compr}

The $\mathrm{Comprehension}$ rule forms $\rbr{f(s_1,\dots,s_n) : s_i \rin S_i}$ for
a function $f\colon S_1\times\dots\times S_n\to S$, and the admissible class of
$f$ is left open although it governs both the strength of the theory and its
safety. The rule is the image of a function over sets already built ---
separation- and replacement-flavoured rather than naive comprehension --- so the
expected outcome is that no Russell set arises; but this must be argued, with
care in the cross-colour case, where $f$ must respect the opacity of atoms. The
cleanest argument is semantic and is already to hand: the $\mathrm{HF}$ model
witnesses consistency, so fixing $f$ to range over the metatheory's set functions
and reading the rule in $\mathcal{M}$ shows the theory paradox-free relative to
$\mathrm{HF}$. Specifying $f$ precisely, and confirming that the cross-colour
instances preserve well-formedness, is the remaining work.

\subsection{Placing the construction}
\label{sec:rigor:place}

Finally, the construction should be located among its neighbours, which would
both guard against reinvention and sharpen its claims. Three comparisons are
immediate. The permutation action and the atoms align it with the
Fraenkel--Mostowski permutation models and their nominal descendants
\cite{DBLP:journals/fac/GabbayP02}; the non-well-founded reading aligns it with
Aczel \cite{aczel1988nwf}; and the colour barrier --- red cannot see into black
--- is a two-level \emph{stratification}, inviting comparison with the typed and
stratified set theories. Most pointedly, red and black \emph{mutually interpret}
one another, so the construction is an instance of mutual interpretation between
set theories; whether the two copies are \emph{bi-interpretable}, and in what
precise sense the universe of set theories promised by the abstract is a
\emph{multiverse}, are exactly the questions that would connect these notes to the
contemporary study of interpretations between set theories. We regard settling
them as the natural completion of the present work.

None of these obligations looks out of reach for the finitary core, where each
reduces to a statement about hereditarily finite sets and several are corollaries
of the model already built. Together they are what would carry the construction
from a suggestive picture to a theorem --- and, not incidentally, they are close
to the list of lemmas a formalisation would have to state in any case.
 (after rset.nwf, before rset.future).
%% rset.colors.tex
%% Red/black colour layer for the two intertwined copies, plus the FM symbols
%% used by rset.mainthm.tex and rset.applications.tex.
%%
%% EVERYTHING routes through \rd (red copy = player) and \bk (black copy =
%% opponent): redefine just those two macros to change the whole scheme (e.g.
%% to make the "black" copy dark blue for legibility).  Every definition is
%% \providecommand-guarded, so it is safe to \input this file from several
%% places and harmless if rset.local already defines these names.
%%
%% Requires \usepackage{xcolor} (already loaded in rset.tex).

% --- the two copies -------------------------------------------------------
\providecommand{\rd}[1]{\textcolor{red}{#1}}     % red copy  (player)
\providecommand{\bk}[1]{\textcolor{black}{#1}}   % black copy (opponent)

% --- theory names and atom supplies --------------------------------------
\providecommand{\rSetX}{\rd{\mathrm{Set}[X]}}    % red   Set[X]
\providecommand{\bSetX}{\bk{\mathrm{Set}[X]}}    % black Set[X]
\providecommand{\rX}{\rd{X}}                      % red   atom supply
\providecommand{\bX}{\bk{X}}                      % black atom supply

% --- empty set ------------------------------------------------------------
\providecommand{\remp}{\rd{\emptyset}}
\providecommand{\bemp}{\bk{\emptyset}}

% --- membership (relations) ----------------------------------------------
\providecommand{\rin}{\mathrel{\rd{\in}}}
\providecommand{\bin}{\mathrel{\bk{\in}}}
\providecommand{\rnin}{\mathrel{\rd{\notin}}}
\providecommand{\bnin}{\mathrel{\bk{\notin}}}

% --- binary operations ----------------------------------------------------
\providecommand{\rcup}{\mathbin{\rd{\cup}}}
\providecommand{\bcup}{\mathbin{\bk{\cup}}}
\providecommand{\rcap}{\mathbin{\rd{\cap}}}
\providecommand{\bcap}{\mathbin{\bk{\cap}}}
\providecommand{\rsm}{\mathbin{\rd{\setminus}}}
\providecommand{\bsm}{\mathbin{\bk{\setminus}}}

% --- comprehension braces of each theory ---------------------------------
\providecommand{\rbr}[1]{\rd{\{}\,#1\,\rd{\}}}
\providecommand{\bbr}[1]{\bk{\{}\,#1\,\bk{\}}}

% --- colour-blind membership (sees through the colour barrier) ------------
\providecommand{\cbin}{\mathrel{\in_{\!\star}}}

% --- FM symbols (shared with rset.stdfmset.tex) --------------------------
\providecommand{\Atoms}{\mathbb{A}}
\providecommand{\Sym}{\mathrm{Sym}}
\providecommand{\pow}{\mathcal{P}}
\providecommand{\supp}{\mathrm{supp}}
\providecommand{\fs}{\mathrm{fs}}
\providecommand{\new}{\mathbin{\reflectbox{\ensuremath{\mathsf{N}}}}}
\providecommand{\atomsof}{\mathrm{atoms}}
\providecommand{\Red}{\rd{\mathsf{R}}}           % the red HF universe
   % red/black colour layer + FM symbols (providecommand-guarded)

\section{Toward a rigorous construction}
\label{sec:rigor}

These are notes, and the constructions above are pitched at the level of intent:
they say what the theory is meant to be and exhibit the model it is meant to
carry, but several of the moves are made by notation where they ought to be made
by theorem. We close by setting down, as a checklist of definitions to fix and
obligations to discharge, what a fully rigorous development would require. We
group them under the headings a careful referee would use --- the fixpoint, the
metatheory, the membership theory, comprehension, and the construction's place
among its neighbours --- and note that for the finitary core each obligation
reduces to a statement about hereditarily finite sets.

\subsection{The knotted fixpoint, via algebraic set theory}
\label{sec:rigor:fixpoint}

The central move, $\rX = \bSetX$ and $\bX = \rSetX$, is at present a definition by
notation; it must be exhibited as a fixed point that provably exists. Algebraic
set theory ($\mathsf{AST}$) is, we think, the right basis for this, because it is
built precisely to produce universes of sets as fixed points of a
power\emph{class} operator while dodging the size obstruction that defeats the
naive power\emph{set} construction \cite{joyal1995algebraic,awodey2008ast}.

Recall its shape. One works in a Heyting category $\mathcal{C}$ of \emph{classes}
with a distinguished collection of \emph{small} maps and a \emph{powerclass}
functor $\mathcal{P}_s\colon\mathcal{C}\to\mathcal{C}$, sending a class to the
class of its \emph{small} sub-objects. The universe of sets is the free
$\mathcal{P}_s$-algebra: an object $U$ with a structural isomorphism
$\mathcal{P}_s U\cong U$, and the central theorem is that this initial
$\mathsf{ZF}$-algebra exists in any category with small maps and models the axioms
of set theory, with membership recovered from the algebra structure. Restricting
$\mathcal{P}_s$ to \emph{small} sub-objects, rather than arbitrary sub-classes, is
exactly what makes $\mathcal{P}_s U\cong U$ possible where Cantor forbids
$\mathcal{P}X\cong X$ --- so the size problem flagged in Section~\ref{sec:nwf} is
dissolved at the level of the framework, not patched case by case.

Two features of $\mathsf{AST}$ fit the present construction almost too well.
First, it handles atoms by a coproduct: the cumulative hierarchy over an object
$A$ of atoms is the initial algebra of $Y\mapsto A + \mathcal{P}_s Y$, in which the
points of $A$ acquire no members --- they are atoms because the free construction
\emph{forgets} whatever structure they carried. That forgetfulness is precisely
our colour opacity (Remark~\ref{rem:group}): red sees the points of its
atom-supply but not their black interior. Second, the universe is constructed as
a $W$-type quotiented by a bounded bisimulation, so the development is
constructive and portable --- and, not incidentally, close to what a
formalisation would have to build.

The knot is then a two-sorted instance. Take the endofunctor on
$\mathcal{C}\times\mathcal{C}$
\[
  H(R,B) \;=\; \bigl(\,B + \mathcal{P}_s R,\ \ R + \mathcal{P}_s B\,\bigr),
\]
whose initial algebra is a pair $(\rX,\bX)$ with $\rX\cong\bX + \mathcal{P}_s\rX$
and $\bX\cong\rX + \mathcal{P}_s\bX$: each colour's universe is its own sets
together with the other colour's universe \emph{as atoms} --- precisely
$\rX=\bSetX$ and $\bX=\rSetX$. The obligation thus becomes to run the
$\mathsf{AST}$ existence and modelling theorems for $H$ in place of
$\mathcal{P}_s$ alone. The ingredients are unchanged --- coproduct, powerclass,
$W$-types, bounded bisimulation --- so one expects the same proof to go through
two-sorted; what it would establish is that the knotted universe \emph{exists} and
that each colour \emph{models} $\mathsf{ZFA}$, discharging in a single stroke both
this obligation and the extensionality and separation obligations of
Section~\ref{sec:rigor:ext}.

Three riders complete the picture. For the \emph{finitary} core one replaces
$\mathcal{P}_s$ by the finite-powerclass functor; being finitary, the initial
algebra of the corresponding $H$ exists already at the level of \emph{sets}, as
the colimit of the initial $\omega$-chain $\varnothing\to H(\varnothing)\to
H^2(\varnothing)\to\cdots$, with no class machinery required --- and it is the
$\mathrm{HF}$ universe of Section~\ref{sec:fmmodel}. For the non-well-founded
reading of Section~\ref{sec:nwf} one takes the \emph{final coalgebra} of $H$ in
place of the initial algebra: at set size for the finitary functor, and in the
class framework --- where it coincides with Aczel's universe --- for the full one.
And a caveat worth stating: the initial $\mathsf{ZF}$-algebra models
\emph{constructive} set theory by default, so recovering classical
$\mathsf{ZFA}$/$\mathsf{FM}$ requires the corresponding classical principles
(power set, full separation, excluded middle) on the class category. With those in
place, $\mathsf{AST}$ supplies exactly the existence theorem these notes have so
far only gestured at, and supplies it in a form that already carries the
membership and extensionality structure the rest of this section asks for.

\subsection{Metatheory and consistency strength}
\label{sec:rigor:meta}

The notes never name the theory in which they are conducted, and the headline
claim --- atoms without infinite risk --- has content only relative to a fixed
metatheory. The obligation is to state that metatheory and to cash the risk claim
as a statement of interpretability. The natural target is hereditarily finite set
theory, equivalently first-order arithmetic: the model $\mathcal{M}$ of
Theorem~\ref{thm:rep} is itself an interpretation of the finitary reflective
theory in $\mathrm{HF}$, so the two are equiconsistent and the ``risk'' is exactly
that of arithmetic --- no infinite set of primitive atoms is ever posited. The
representation theorem, in other words, already \emph{is} a relative-consistency
proof; what remains is to state it as one. The infinitary extension of
Section~\ref{sec:fmmodel:infinitary} raises the target by a single $\omega$: it is
interpretable in $\mathsf{Z}$ --- or in $\mathsf{ZF}$ with foundation reinstated
per colour --- the lone infinity being black infinity. Making ``finite versus
infinite risk'' precise in this way is, we expect, the first thing a set theorist
would ask for, and it is the claim on which the paper's title rests.

\subsection{Extensionality and the membership theory}
\label{sec:rigor:ext}

Extensionality is not visibly secured by the equations of the formal theory, which
give reflexivity, symmetry, transitivity, and idempotence of $\rcup$ and $\rcap$,
but not that sets with the same members are equal. A rigorous development must
either adopt extensionality as an axiom or derive it, and must then prove that
provable red equality coincides with sameness of members --- the property the
proof of Theorem~\ref{thm:rep} leans on. The membership rules need a matching
completeness theorem. As presented they decide $\rin$ on unions, intersections,
differences, and singletons, but no rule governs membership in a nested
$\rbr{S}$, and one wants a proof that for $\mathrm{HF}$ sets every true $e\rin S$
and $e\rnin S$ is derivable and that the relation is decidable. Until then, the
remark that the syntactic theory is ``too fine grained'' and is corrected by
equations is a programme rather than a theorem. We note that the
algebraic-set-theory route of Section~\ref{sec:rigor:fixpoint} would discharge
this obligation along with existence: its modelling theorem delivers
extensionality for the constructed universe directly, so on that route the
equational theory need only be shown \emph{sound} for the model rather than
independently complete.

\subsection{Comprehension and freedom from paradox}
\label{sec:rigor:compr}

The $\mathrm{Comprehension}$ rule forms $\rbr{f(s_1,\dots,s_n) : s_i \rin S_i}$ for
a function $f\colon S_1\times\dots\times S_n\to S$, and the admissible class of
$f$ is left open although it governs both the strength of the theory and its
safety. The rule is the image of a function over sets already built ---
separation- and replacement-flavoured rather than naive comprehension --- so the
expected outcome is that no Russell set arises; but this must be argued, with
care in the cross-colour case, where $f$ must respect the opacity of atoms. The
cleanest argument is semantic and is already to hand: the $\mathrm{HF}$ model
witnesses consistency, so fixing $f$ to range over the metatheory's set functions
and reading the rule in $\mathcal{M}$ shows the theory paradox-free relative to
$\mathrm{HF}$. Specifying $f$ precisely, and confirming that the cross-colour
instances preserve well-formedness, is the remaining work.

\subsection{Placing the construction}
\label{sec:rigor:place}

Finally, the construction should be located among its neighbours, which would
both guard against reinvention and sharpen its claims. Three comparisons are
immediate. The permutation action and the atoms align it with the
Fraenkel--Mostowski permutation models and their nominal descendants
\cite{DBLP:journals/fac/GabbayP02}; the non-well-founded reading aligns it with
Aczel \cite{aczel1988nwf}; and the colour barrier --- red cannot see into black
--- is a two-level \emph{stratification}, inviting comparison with the typed and
stratified set theories. Most pointedly, red and black \emph{mutually interpret}
one another, so the construction is an instance of mutual interpretation between
set theories; whether the two copies are \emph{bi-interpretable}, and in what
precise sense the universe of set theories promised by the abstract is a
\emph{multiverse}, are exactly the questions that would connect these notes to the
contemporary study of interpretations between set theories. We regard settling
them as the natural completion of the present work.

None of these obligations looks out of reach for the finitary core, where each
reduces to a statement about hereditarily finite sets and several are corollaries
of the model already built. Together they are what would carry the construction
from a suggestive picture to a theorem --- and, not incidentally, they are close
to the list of lemmas a formalisation would have to state in any case.
 (after rset.nwf, before rset.future).
%% rset.colors.tex
%% Red/black colour layer for the two intertwined copies, plus the FM symbols
%% used by rset.mainthm.tex and rset.applications.tex.
%%
%% EVERYTHING routes through \rd (red copy = player) and \bk (black copy =
%% opponent): redefine just those two macros to change the whole scheme (e.g.
%% to make the "black" copy dark blue for legibility).  Every definition is
%% \providecommand-guarded, so it is safe to \input this file from several
%% places and harmless if rset.local already defines these names.
%%
%% Requires \usepackage{xcolor} (already loaded in rset.tex).

% --- the two copies -------------------------------------------------------
\providecommand{\rd}[1]{\textcolor{red}{#1}}     % red copy  (player)
\providecommand{\bk}[1]{\textcolor{black}{#1}}   % black copy (opponent)

% --- theory names and atom supplies --------------------------------------
\providecommand{\rSetX}{\rd{\mathrm{Set}[X]}}    % red   Set[X]
\providecommand{\bSetX}{\bk{\mathrm{Set}[X]}}    % black Set[X]
\providecommand{\rX}{\rd{X}}                      % red   atom supply
\providecommand{\bX}{\bk{X}}                      % black atom supply

% --- empty set ------------------------------------------------------------
\providecommand{\remp}{\rd{\emptyset}}
\providecommand{\bemp}{\bk{\emptyset}}

% --- membership (relations) ----------------------------------------------
\providecommand{\rin}{\mathrel{\rd{\in}}}
\providecommand{\bin}{\mathrel{\bk{\in}}}
\providecommand{\rnin}{\mathrel{\rd{\notin}}}
\providecommand{\bnin}{\mathrel{\bk{\notin}}}

% --- binary operations ----------------------------------------------------
\providecommand{\rcup}{\mathbin{\rd{\cup}}}
\providecommand{\bcup}{\mathbin{\bk{\cup}}}
\providecommand{\rcap}{\mathbin{\rd{\cap}}}
\providecommand{\bcap}{\mathbin{\bk{\cap}}}
\providecommand{\rsm}{\mathbin{\rd{\setminus}}}
\providecommand{\bsm}{\mathbin{\bk{\setminus}}}

% --- comprehension braces of each theory ---------------------------------
\providecommand{\rbr}[1]{\rd{\{}\,#1\,\rd{\}}}
\providecommand{\bbr}[1]{\bk{\{}\,#1\,\bk{\}}}

% --- colour-blind membership (sees through the colour barrier) ------------
\providecommand{\cbin}{\mathrel{\in_{\!\star}}}

% --- FM symbols (shared with rset.stdfmset.tex) --------------------------
\providecommand{\Atoms}{\mathbb{A}}
\providecommand{\Sym}{\mathrm{Sym}}
\providecommand{\pow}{\mathcal{P}}
\providecommand{\supp}{\mathrm{supp}}
\providecommand{\fs}{\mathrm{fs}}
\providecommand{\new}{\mathbin{\reflectbox{\ensuremath{\mathsf{N}}}}}
\providecommand{\atomsof}{\mathrm{atoms}}
\providecommand{\Red}{\rd{\mathsf{R}}}           % the red HF universe
   % red/black colour layer + FM symbols (providecommand-guarded)

\section{Toward a rigorous construction}
\label{sec:rigor}

These are notes, and the constructions above are pitched at the level of intent:
they say what the theory is meant to be and exhibit the model it is meant to
carry, but several of the moves are made by notation where they ought to be made
by theorem. We close by setting down, as a checklist of definitions to fix and
obligations to discharge, what a fully rigorous development would require. We
group them under the headings a careful referee would use --- the fixpoint, the
metatheory, the membership theory, comprehension, and the construction's place
among its neighbours --- and note that for the finitary core each obligation
reduces to a statement about hereditarily finite sets.

\subsection{The knotted fixpoint, via algebraic set theory}
\label{sec:rigor:fixpoint}

The central move, $\rX = \bSetX$ and $\bX = \rSetX$, is at present a definition by
notation; it must be exhibited as a fixed point that provably exists. Algebraic
set theory ($\mathsf{AST}$) is, we think, the right basis for this, because it is
built precisely to produce universes of sets as fixed points of a
power\emph{class} operator while dodging the size obstruction that defeats the
naive power\emph{set} construction \cite{joyal1995algebraic,awodey2008ast}.

Recall its shape. One works in a Heyting category $\mathcal{C}$ of \emph{classes}
with a distinguished collection of \emph{small} maps and a \emph{powerclass}
functor $\mathcal{P}_s\colon\mathcal{C}\to\mathcal{C}$, sending a class to the
class of its \emph{small} sub-objects. The universe of sets is the free
$\mathcal{P}_s$-algebra: an object $U$ with a structural isomorphism
$\mathcal{P}_s U\cong U$, and the central theorem is that this initial
$\mathsf{ZF}$-algebra exists in any category with small maps and models the axioms
of set theory, with membership recovered from the algebra structure. Restricting
$\mathcal{P}_s$ to \emph{small} sub-objects, rather than arbitrary sub-classes, is
exactly what makes $\mathcal{P}_s U\cong U$ possible where Cantor forbids
$\mathcal{P}X\cong X$ --- so the size problem flagged in Section~\ref{sec:nwf} is
dissolved at the level of the framework, not patched case by case.

Two features of $\mathsf{AST}$ fit the present construction almost too well.
First, it handles atoms by a coproduct: the cumulative hierarchy over an object
$A$ of atoms is the initial algebra of $Y\mapsto A + \mathcal{P}_s Y$, in which the
points of $A$ acquire no members --- they are atoms because the free construction
\emph{forgets} whatever structure they carried. That forgetfulness is precisely
our colour opacity (Remark~\ref{rem:group}): red sees the points of its
atom-supply but not their black interior. Second, the universe is constructed as
a $W$-type quotiented by a bounded bisimulation, so the development is
constructive and portable --- and, not incidentally, close to what a
formalisation would have to build.

The knot is then a two-sorted instance. Take the endofunctor on
$\mathcal{C}\times\mathcal{C}$
\[
  H(R,B) \;=\; \bigl(\,B + \mathcal{P}_s R,\ \ R + \mathcal{P}_s B\,\bigr),
\]
whose initial algebra is a pair $(\rX,\bX)$ with $\rX\cong\bX + \mathcal{P}_s\rX$
and $\bX\cong\rX + \mathcal{P}_s\bX$: each colour's universe is its own sets
together with the other colour's universe \emph{as atoms} --- precisely
$\rX=\bSetX$ and $\bX=\rSetX$. The obligation thus becomes to run the
$\mathsf{AST}$ existence and modelling theorems for $H$ in place of
$\mathcal{P}_s$ alone. The ingredients are unchanged --- coproduct, powerclass,
$W$-types, bounded bisimulation --- so one expects the same proof to go through
two-sorted; what it would establish is that the knotted universe \emph{exists} and
that each colour \emph{models} $\mathsf{ZFA}$, discharging in a single stroke both
this obligation and the extensionality and separation obligations of
Section~\ref{sec:rigor:ext}.

Three riders complete the picture. For the \emph{finitary} core one replaces
$\mathcal{P}_s$ by the finite-powerclass functor; being finitary, the initial
algebra of the corresponding $H$ exists already at the level of \emph{sets}, as
the colimit of the initial $\omega$-chain $\varnothing\to H(\varnothing)\to
H^2(\varnothing)\to\cdots$, with no class machinery required --- and it is the
$\mathrm{HF}$ universe of Section~\ref{sec:fmmodel}. For the non-well-founded
reading of Section~\ref{sec:nwf} one takes the \emph{final coalgebra} of $H$ in
place of the initial algebra: at set size for the finitary functor, and in the
class framework --- where it coincides with Aczel's universe --- for the full one.
And a caveat worth stating: the initial $\mathsf{ZF}$-algebra models
\emph{constructive} set theory by default, so recovering classical
$\mathsf{ZFA}$/$\mathsf{FM}$ requires the corresponding classical principles
(power set, full separation, excluded middle) on the class category. With those in
place, $\mathsf{AST}$ supplies exactly the existence theorem these notes have so
far only gestured at, and supplies it in a form that already carries the
membership and extensionality structure the rest of this section asks for.

\subsection{Metatheory and consistency strength}
\label{sec:rigor:meta}

The notes never name the theory in which they are conducted, and the headline
claim --- atoms without infinite risk --- has content only relative to a fixed
metatheory. The obligation is to state that metatheory and to cash the risk claim
as a statement of interpretability. The natural target is hereditarily finite set
theory, equivalently first-order arithmetic: the model $\mathcal{M}$ of
Theorem~\ref{thm:rep} is itself an interpretation of the finitary reflective
theory in $\mathrm{HF}$, so the two are equiconsistent and the ``risk'' is exactly
that of arithmetic --- no infinite set of primitive atoms is ever posited. The
representation theorem, in other words, already \emph{is} a relative-consistency
proof; what remains is to state it as one. The infinitary extension of
Section~\ref{sec:fmmodel:infinitary} raises the target by a single $\omega$: it is
interpretable in $\mathsf{Z}$ --- or in $\mathsf{ZF}$ with foundation reinstated
per colour --- the lone infinity being black infinity. Making ``finite versus
infinite risk'' precise in this way is, we expect, the first thing a set theorist
would ask for, and it is the claim on which the paper's title rests.

\subsection{Extensionality and the membership theory}
\label{sec:rigor:ext}

Extensionality is not visibly secured by the equations of the formal theory, which
give reflexivity, symmetry, transitivity, and idempotence of $\rcup$ and $\rcap$,
but not that sets with the same members are equal. A rigorous development must
either adopt extensionality as an axiom or derive it, and must then prove that
provable red equality coincides with sameness of members --- the property the
proof of Theorem~\ref{thm:rep} leans on. The membership rules need a matching
completeness theorem. As presented they decide $\rin$ on unions, intersections,
differences, and singletons, but no rule governs membership in a nested
$\rbr{S}$, and one wants a proof that for $\mathrm{HF}$ sets every true $e\rin S$
and $e\rnin S$ is derivable and that the relation is decidable. Until then, the
remark that the syntactic theory is ``too fine grained'' and is corrected by
equations is a programme rather than a theorem. We note that the
algebraic-set-theory route of Section~\ref{sec:rigor:fixpoint} would discharge
this obligation along with existence: its modelling theorem delivers
extensionality for the constructed universe directly, so on that route the
equational theory need only be shown \emph{sound} for the model rather than
independently complete.

\subsection{Comprehension and freedom from paradox}
\label{sec:rigor:compr}

The $\mathrm{Comprehension}$ rule forms $\rbr{f(s_1,\dots,s_n) : s_i \rin S_i}$ for
a function $f\colon S_1\times\dots\times S_n\to S$, and the admissible class of
$f$ is left open although it governs both the strength of the theory and its
safety. The rule is the image of a function over sets already built ---
separation- and replacement-flavoured rather than naive comprehension --- so the
expected outcome is that no Russell set arises; but this must be argued, with
care in the cross-colour case, where $f$ must respect the opacity of atoms. The
cleanest argument is semantic and is already to hand: the $\mathrm{HF}$ model
witnesses consistency, so fixing $f$ to range over the metatheory's set functions
and reading the rule in $\mathcal{M}$ shows the theory paradox-free relative to
$\mathrm{HF}$. Specifying $f$ precisely, and confirming that the cross-colour
instances preserve well-formedness, is the remaining work.

\subsection{Placing the construction}
\label{sec:rigor:place}

Finally, the construction should be located among its neighbours, which would
both guard against reinvention and sharpen its claims. Three comparisons are
immediate. The permutation action and the atoms align it with the
Fraenkel--Mostowski permutation models and their nominal descendants
\cite{DBLP:journals/fac/GabbayP02}; the non-well-founded reading aligns it with
Aczel \cite{aczel1988nwf}; and the colour barrier --- red cannot see into black
--- is a two-level \emph{stratification}, inviting comparison with the typed and
stratified set theories. Most pointedly, red and black \emph{mutually interpret}
one another, so the construction is an instance of mutual interpretation between
set theories; whether the two copies are \emph{bi-interpretable}, and in what
precise sense the universe of set theories promised by the abstract is a
\emph{multiverse}, are exactly the questions that would connect these notes to the
contemporary study of interpretations between set theories. We regard settling
them as the natural completion of the present work.

None of these obligations looks out of reach for the finitary core, where each
reduces to a statement about hereditarily finite sets and several are corollaries
of the model already built. Together they are what would carry the construction
from a suggestive picture to a theorem --- and, not incidentally, they are close
to the list of lemmas a formalisation would have to state in any case.
